\subsection{REPRESENTATIONS ON HOMOMORPHISM MODULES}

Again let \(\mathfrak{g}_1\) and \(\mathfrak{g}_2\) be two Lie algebras over K and \(\mathbf{M}_i\) a \(\mathfrak{g}_i\) -module ( \(i = 1, 2\) ). Let \(\mathrm{U}_i\) be the enveloping algebra of \(\mathfrak{g}_i\) and \(\sigma_i\) the canonical mapping of \(\mathfrak{g}_i\) into \(\mathrm{U}_i\) . Then \(\mathbf{M}_i\) is a left \(\mathrm{U}_i\) -module and hence \(\mathcal{L}_{\mathrm{K}}(\mathrm{M}_1, \mathrm{M}_2)\) has a canonical left \((\mathrm{U}_1^0 \otimes \mathrm{U}_2)\) -module structure. Now \(\mathrm{U}_1^0 \otimes_{\mathrm{K}} \mathrm{U}_2\) is the enveloping algebra of \(\mathfrak{g}_1^0 \times \mathfrak{g}_2\) and the mapping

\[
(x _ {1}, x _ {2}) \mapsto \sigma_ {1} (x _ {1}) \otimes 1 + 1 \otimes \sigma_ {2} (x _ {2})
\]

is the canonical mapping of \(\mathfrak{g}_1^0\times \mathfrak{g}_2\) into this enveloping algebra. Hence there exists a \((\mathfrak{g}_1^0\times \mathfrak{g}_2)\) -module structure on \(\mathbf{M} = \mathcal{L}_{\mathbb{K}}(\mathbf{M}_1,\mathbf{M}_2)\) such that

\[
\begin{array}{r l} ((x _ {1}, x _ {2}) _ {\mathrm{M}}. u) \cdot m _ {1} & = ((\sigma_ {1} (x _ {1}) \otimes 1 + 1 \otimes \sigma_ {2} (x _ {2})) \cdot u) \cdot m _ {1} \\ & = u ((x _ {1}) _ {\mathrm{M} _ {1}}. m _ {1}) + (x _ {2}) _ {\mathrm{M} _ {2}}. u (m _ {1}) \end{array}\tag{4}
\]

for all \(u \in \mathcal{L}_{\mathbb{K}}(\mathrm{M}_1, \mathrm{M}_2)\) , \(m_1 \in \mathrm{M}_1\) . This structure defines a representation of \(\mathfrak{g}_1^0 \times \mathfrak{g}_2\) on \(\mathbf{M}\) .

If now \(\mathfrak{g}_1 = \mathfrak{g}_2 = \mathfrak{g}\) , the homomorphism \(x \mapsto (-x, x)\) of \(\mathfrak{g}\) into \(\mathfrak{g}^0 \times \mathfrak{g}\) , composed with the above representation, defines a representation of g on M and hence a g-module structure on M such that

\[
(x _ {\mathrm{M}}. u). m _ {1} = x _ {\mathrm{M} _ {2}}. u (m _ {1}) - u (x _ {\mathrm{M} _ {1}}. m _ {1})\tag{5}
\]

or

\[
x _ {\mathrm{M}}. u = x _ {\mathrm{M} _ {2}} u - u x _ {\mathrm{M} _ {1}}.\tag{6}
\]

Combining these results with Proposition 2, we see that:

PROPOSITION 3. Let \(\mathfrak{g}\) be a Lie algebra over \(\mathbf{K}\) and \(\mathbf{M}_i\) a \(\mathfrak{g}\) -module ( \(1 \leqslant i \leqslant n + 1\) ).

Let \(\mathbf{N}\) be the K-module \(\mathcal{L}_{\mathbf{K}}(\mathbf{M}_1, \ldots, \mathbf{M}_n; \mathbf{M}_{n+1})\) of multilinear mappings of \(\prod_{i=1}^{n} \mathbf{M}_i\) into \(\mathbf{M}_{n+1}\) . There exists one and only one g-module structure on \(\mathbf{N}\) such that

\[
\begin{array}{r l} (x _ {N}. u) (m _ {1}, \ldots , m _ {n}) & = - \sum_ {i = 1} ^ {n} u (m _ {1}, \ldots , x _ {M _ {i}}. m _ {i}, \ldots , m _ {n}) \\ & \quad + x _ {M _ {n + 1}}. u (m _ {1}, \ldots , m _ {n}) \end{array}\tag{7}
\]

for all \(x \in \mathfrak{g}\) , \(u \in \mathbf{N}\) and \(m_i \in \mathbf{M}_i\) ( \(1 \leqslant i \leqslant n\) ).

In particular, let g be a Lie algebra over K and M a g-module and consider K as a trivial g-module. Proposition 3 defines a g-module structure on \(\mathcal{L}_{\mathrm{K}}(\mathrm{M},\mathrm{K}) = \mathrm{M}^{*}\) . The corresponding representation is called the dual representation of the representation \(x\mapsto x_{\mathrm{M}}\) . We have:

\[
(x _ {\mathrm{M} ^ {*}}. f) (m) = - f (x _ {\mathrm{M}}. m)\tag{8}
\]

for all \(x \in \mathfrak{g}, f \in \mathbf{M}^*\) , \(m \in \mathbf{M}\) . In other words:

\[
x _ {\mathrm{M} ^ {*}} = - ^ {t} x _ {\mathrm{M}}.\tag{9}
\]

When K is a field and M is finite-dimensional, the g-module M is simple (resp. semi-simple) if and only if the g-module M* is simple (resp. semi-simple).

PROPOSITION 4. Let \(\mathbf{M}_1, \mathbf{M}_2\) be two g-modules. The canonical K-linear mappings (Algebra, Chapter II, § 4, no. 2, Proposition 2 and no. 1, Proposition 1):

\[
\mathrm{M} _ {1} ^ {*} \otimes_ {\mathrm{K}} \mathrm{M} _ {2} \xrightarrow {\Phi} \mathscr {L} _ {\mathrm{K}} (\mathrm{M} _ {1}, \mathrm{M} _ {2}), \quad \mathscr {L} _ {\mathrm{K}} (\mathrm{M} _ {1}, \mathrm{M} _ {2} ^ {*}) \xrightarrow {\Psi} (\mathrm{M} _ {1} \otimes_ {\mathrm{K}} \mathrm{M} _ {2}) ^ {*}
\]

(where the second is bijective) are \(\mathfrak{g}\) -module homomorphisms.

We write

\[
\mathrm{N} = \mathrm{M} _ {1} ^ {*} \otimes \mathrm{M} _ {2}, \quad \mathrm{P} = \mathscr {L} (\mathrm{M} _ {1}, \mathrm{M} _ {2}), \quad \mathrm{Q} = \mathscr {L} (\mathrm{M} _ {1}, \mathrm{M} _ {2} ^ {*}), \quad \mathrm{R} = (\mathrm{M} _ {1} \otimes \mathrm{M} _ {2}) ^ {*}.
\]

Then, for \(x \in \mathfrak{g}, f \in \mathbf{M}_1^*\) , \(m_1 \in \mathbf{M}_1\) , \(m_2 \in \mathbf{M}_2\) ,

\[
\begin{array}{r l} & {((\phi x _ {\mathrm{N}}) (f \otimes m _ {2})). m _ {1} = (\phi (x _ {\mathrm{M} _ {1} ^ {*}} f \otimes m _ {2} + f \otimes x _ {\mathrm{M} _ {2}} m _ {2})). m _ {1}} \\ & {\qquad = \langle x _ {\mathrm{M} _ {1} ^ {*}} f, m _ {1} \rangle m _ {2} + \langle f, m _ {1} \rangle x _ {\mathrm{M} _ {2}} m _ {2}} \\ & {((x _ {\mathrm{P}} \phi) (f \otimes m _ {2})). m _ {1} = x _ {\mathrm{M} _ {2}} (\phi (f \otimes m _ {2}). m _ {1}) - \phi (f \otimes m _ {2}) (x _ {\mathrm{M} _ {1}} m _ {1})} \\ & {\qquad = \langle f, m _ {1} \rangle x _ {\mathrm{M} _ {2}} m _ {2} - \langle f, x _ {\mathrm{M} _ {1}} m _ {1} \rangle m _ {2}} \end{array}
\]

and hence \(\phi x_{\mathbf{N}} = x_{\mathbb{P}}\phi\) . On the other hand, for \(x \in \mathfrak{g}\) , \(u \in \mathcal{L}(\mathrm{M}_1, \mathrm{M}_2^*)\) , \(m_1 \in \mathrm{M}_1\) , \(m_2 \in \mathrm{M}_2\) :

\[
\begin{array}{r l} (\psi x _ {\mathbf {Q}} u) (m _ {1} \otimes m _ {2}) & = \langle (x _ {\mathbf {Q}} u). m _ {1}, m _ {2} \rangle = \langle x _ {\mathrm{M} _ {2} ^ {*}} u m _ {1} - u x _ {\mathrm{M} _ {1}} m _ {1}, m _ {2} \rangle \\ (x _ {\mathrm{R}} \psi u) (m _ {1} \otimes m _ {2}) & = - \langle \psi u, x _ {\mathrm{M} _ {1}} m _ {1} \otimes m _ {2} + m _ {1} \otimes x _ {\mathrm{M} _ {2}} m _ {2} \rangle \\ & = - \langle u x _ {\mathrm{M} _ {1}} m _ {1}, m _ {2} \rangle - \langle u m _ {1}, x _ {\mathrm{M} _ {2}} m _ {2} \rangle \end{array}
\]

and hence \(\psi x_{Q} = x_{R}\psi\) , which completes the proof.

The g-modules \(\mathcal{L}(\mathrm{M}_1,\mathrm{M}_2^*)\) and \((\mathrm{M}_1\otimes \mathrm{M}_2)^*\) are identified under the isomorphism \(\psi\) . If \(\mathbf{M}_1\) and \(\mathbf{M}_2\) have finite bases, \(\phi\) is an isomorphism (Algebra, Chapter II, § 4, no. 2, Proposition 2), which allows us to identify the g-modules \(\mathbf{M}_1^*\otimes \mathbf{M}_2\) and \(\mathcal{L}(\mathrm{M}_1,\mathrm{M}_2)\) ; in that case, we can therefore identify the g-modules \(\mathbf{M}_1^*\otimes \mathbf{M}_2^*\) , \(\mathcal{L}(\mathrm{M}_1,\mathbf{M}_2^*)\) and \((\mathrm{M}_1\otimes \mathrm{M}_2)^*\) .

